{\light\fontsize{11.5}{16}\selectfont\color{SpiritNavy}
Linear algebra has a lot of vocabulary, and many of the words sound alike. This glossary gives each term
from MATH 1850U in plain language first, then the precise definition, a small example, and a mistake
students often make with it.\par}

\vspace{14pt}
\lbl[SimcoeBlue]{Anatomy of an entry}\par\vspace{3pt}{\color{FutureBlue}\rule{\linewidth}{0.8pt}}\par\vspace{12pt}

% ---- annotated sample entry (not linked, not bookmarked)
\newcommand{\callout}[2]{\makebox[0pt][r]{\tikz[baseline=-0.6ex]{\fill[Tangerine] (0,0) circle (6.5pt);
  \node[white,font=\bfseries\fontsize{8}{8}\selectfont] at (0,0) {#1};}\hspace{6pt}}#2}
\noindent\begin{minipage}[t]{\leftcol}\raggedright
  {\bfseries\fontsize{12.6}{14.6}\selectfont\color{FutureBlue}\callout{1}{Span}\par}%
\end{minipage}\hfill
\begin{minipage}[t]{\dimexpr\linewidth-\leftcol-0.22in\relax}\raggedright\fontsize{9.6}{13.4}\selectfont
  \plainwords{\callout{2}{Everything you can build from a set of vectors using linear combinations.}}
  \begin{defbox}\callout{3}{\lbl[SimcoeBlue]{Definition}}\par\vspace{2pt}
    The set of all \textcolor{SimcoeBlue}{\medium linear combinations} $k_1\mathbf{v}_1 + \dots + k_r\mathbf{v}_r$
    of the vectors in $S$, written $\operatorname{span}(S)$.\end{defbox}
  \begin{exbox}\callout{4}{\lbl[DarkGrey]{Example}}\par\vspace{2pt}
    $\operatorname{span}\{(1, 0, 0), (0, 1, 0)\}$ is the $xy$-plane in $\mathbb{R}^3$.\end{exbox}
  \begin{watchbox}\callout{5}{\lbl[FutureBlue]{Don't confuse with}}\par\vspace{2pt}
    \textcolor{SimcoeBlue}{\medium Basis}. A spanning set can contain redundant vectors, while a basis cannot.\end{watchbox}
  \par\vspace{6pt}{\fontsize{8.6}{12}\selectfont\callout{6}{\lbl[DarkGrey]{See also}}\hspace{5pt}%
    \textcolor{SimcoeBlue}{\medium Spanning set, Basis, Column space}\par}
\end{minipage}

\vspace{14pt}
{\fontsize{9}{12.4}\selectfont
\begin{multicols}{2}\raggedright
\newcommand{\key}[2]{\noindent\makebox[0.3in][l]{\tikz[baseline=-0.6ex]{\fill[Tangerine] (0,0) circle (6.5pt);
  \node[white,font=\bfseries\fontsize{8}{8}\selectfont] at (0,0) {#1};}}\parbox[t]{\dimexpr\linewidth-0.3in}{#2}\par\vspace{5pt}}
\key{1}{\textbf{The term}, as it is written in the course notes. An \textbf{Application} badge under it marks
  terms that come from an applied example rather than the core theory.}
\key{2}{\textbf{In plain words.} Start here if the formal definition does not make sense yet.}
\key{3}{\textbf{Definition.} The precise version to use on assignments and tests.}
\key{4}{\textbf{Example.} A small, concrete case.}
\key{5}{\textbf{Don't confuse with.} A common mix-up for this term and how to tell the two apart.}
\key{6}{\textbf{See also} and \textbf{also called}. Related terms and other names you may see in other
  textbooks or videos.}
\end{multicols}}

\vspace{4pt}
\lbl[SimcoeBlue]{Finding your way around}\par\vspace{3pt}{\color{FutureBlue}\rule{\linewidth}{0.8pt}}\par\vspace{10pt}
{\fontsize{9}{12.4}\selectfont
\begin{itemize}[leftmargin=1em,itemsep=2pt]
  \item[\textcolor{Tangerine}{\rule{4pt}{4pt}}] \textbf{Everything in \textcolor{SimcoeBlue}{blue} is a link.}
    Click a term to jump to its entry.
  \item[\textcolor{Tangerine}{\rule{4pt}{4pt}}] The \textbf{footer} on every page links back to the
    \hyperref[p:contents]{\textcolor{SimcoeBlue}{\medium Contents}},
    \hyperref[p:notation]{\textcolor{SimcoeBlue}{\medium Notation decoder}},
    \hyperref[p:pairs]{\textcolor{SimcoeBlue}{\medium Common mix-ups}} and
    \hyperref[p:index]{\textcolor{SimcoeBlue}{\medium A-Z index}}.
  \item[\textcolor{Tangerine}{\rule{4pt}{4pt}}] Open your PDF viewer's \textbf{bookmarks} panel to see
    every section and term in a sidebar.
  \item[\textcolor{Tangerine}{\rule{4pt}{4pt}}] Terms are grouped in \textbf{course order}, so each section
    builds on the ones before it. Use the A-Z index if you only know the word.
\end{itemize}}
